\section{Capacity on a curve}
\label{sec:ch25-curves}

The capacity of a tube over a curve skeleton is a weighted graph
problem.  We prove the pushforward identity directly from model
integration by parts.  This keeps both the residue degrees and the
nonreduced SNC refinements that occur over the original field visible.

Retain the original BFJ hypotheses \eqref{eq:ch20-standing}, with
\(-\log|\pi|=1\).  Thus \(K\) is complete and discretely valued,
\(k=\widetilde K\) has characteristic zero, and \(C\) is a smooth
connected projective \(K\)-curve.  The reference \(\theta\) is an
arbitrary real semipositive closed model form with ample generic class;
it need not come from one line bundle.  Suppose it is determined on
a given semistable model \(\mathcal C\), with smooth components as in
the regular semistable convention of
\cref{def:ch08-semistable-curve-model}, but without requiring its nodes
to be \(k\)-rational.  A node of thickness greater than one may first
be resolved into its regular chain: this subdivides its edge into
unit-length edges over its node residue field, without changing the
underlying skeleton or retraction.  This is a resolution of the given
model, not a semistable domination assertion for arbitrary models.
If the model does not yet satisfy the irreducible-strata indexing
convention for SNC models, make a fixed node subdivision as in
\cref{prop:snc-cofinal}.  It leaves the underlying graph and
retraction unchanged, although the inserted components can have
multiplicity greater than one.  It is this SNC presentation that
is used when applying \cref{def:ch21-bfj-psh} to the fixed graph.

Write \(\Gamma\) for that skeleton,
\(p:C^{\An}\to\Gamma\) for its retraction and
\(\iota:\Gamma\hookrightarrow C^{\An}\) for its inclusion.  Let
\(\mu_0=\MA_\theta(0)\), viewed on \(\Gamma\).  Its mass is
\(V=(\{\theta\})>0\), and its vertex masses are the degrees of the
determining real class over \(k\).  If \(e\) corresponds to a node
\(z_e\), set
\[
 c_e=[\kappa(z_e):k].
\]
Each edge has its analytic modulus coordinate \(s\).  The same
\(c_e\) is used at both ends and on every subdivision of \(e\).
For a continuous piecewise affine \(T\), our convention is
\[
 \Delta_\Gamma T
 =\sum_x\left(\sum_{e\ni x}c_e\,\partial_eT(x)\right)\delta_x,
 \qquad
 \int_\Gamma F\,\rd\Delta_\Gamma T
 =-\sum_e c_e\int_e F'T'\,\rd s
\]
when \(F\) is also piecewise affine, after a common subdivision.
This is the weighted Laplacian of
\eqref{eq:ch31-graph-laplacian}, not the unit-weight formula
\eqref{eq:graph-laplacian}.  For a continuous \(F\), the assertion
\(\Delta_\Gamma F=\nu\), with \(\nu\) a finite signed measure, means
\(\int F\,\rd\Delta_\Gamma T=\int T\,\rd\nu\) for every continuous
piecewise affine \(T\).

We use the real-form graph lemma
\cref{lem:ch26-weighted-graph-real-bfj}.  Its proof uses the algebraic
curve intersection formula, model integration by parts, uniform
stability and continuity of the bounded operator, and the elementary
weighted Green kernel.  It uses no capacity theorem and no
finite-energy result; this forward reference is therefore independent
of the capacity arguments of this chapter.

Here is its model-level input in the notation needed below.  If \(T\)
is real piecewise affine with rational breakpoints in the original
modulus coordinates, then \(t=T\circ p\) is a real model function and
\begin{equation}
\label{eq:ch25-curve-ma}
 \MA_\theta(t)=\iota_*(\mu_0+\Delta_\Gamma T).
\end{equation}
This is an equality of signed model measures if \(t\) is not psh.
In particular \(\ddc t=\iota_*\Delta_\Gamma T\).  To recall why the
statement holds over \(K\), successive blowups of closed nodes insert
the desired rational points.  The resulting regular SNC model may
have special fiber \(\sum_i b_iE_i\) with \(b_i>1\).
For prescribed real vertex values \(a_i\), the vertical real divisor
is \(D=\sum_i b_i a_iE_i\), not \(\sum_i a_iE_i\).
Adjacent vertices obtained by these node blowups have modulus
distance \(1/(b_i b_j)\), and their intersection has degree \(c_e\).
Consequently the first identity in
\eqref{eq:ch20-curve-laplacian} gives
\[
 b_i(D\cdot E_i)
   =\sum_{j\ne i}b_i b_j(E_j\cdot E_i)(a_j-a_i)
   =\sum_{e\ni i}c_e\,\partial_eT(i).
\]
The left side is exactly the atom in
\eqref{eq:model-mixed-measure}; the reference has degree zero on
new exceptional components by the projection formula.  This proves
\eqref{eq:ch25-curve-ma}.  For example, blowing up a node of a unit
edge inserts its midpoint with multiplicity two.  There is no
requirement that this SNC model have reduced special fiber, and no
claim that a real \(T\) be a rational model function.

\begin{proposition}
\label{prop:ch25-curve-pushforward}
For every bounded \(\theta\)-psh function \(u\) on \(C^{\An}\), the
restriction \(f=u|_\Gamma\) is continuous,
\[
 \mu_0+\Delta_\Gamma f=p_*\MA_\theta(u)\geq0,
 \qquad u\circ p\in\PSH(C,\theta)\cap C^0(C^{\An}),
\]
and
\begin{equation}
\label{eq:ch25-curve-pushforward}
 \iota_*p_*\MA_\theta(u)=\MA_\theta(u\circ p).
\end{equation}
Thus \(\MA_\theta(u)(p^{-1}(A))
=\MA_\theta(u\circ p)(\iota(A))\) for every Borel
\(A\subseteq\Gamma\).  In particular
\eqref{eq:ch25-curve-ma} also holds with \(t=u\circ p\) and
\(T=u|_\Gamma\), without a model-function hypothesis.
\end{proposition}

\begin{proof}
\emph{Step 1: test on the original skeleton.}
Let \(\psi\) be any real \(\theta\)-psh model function; its determining
model is not assumed semistable.  For a real piecewise affine \(T\)
with rational breakpoints, \(t=T\circ p\) is the model function just
constructed.  Model integration by parts
\cref{prop:model-integration-by-parts}, applied on a common model
determining \(\psi,t,\theta\), gives
\begin{equation}
\label{eq:ch25-weighted-skeletal-test}
 \begin{aligned}
 \int_\Gamma T\,\rd(p_*\MA_\theta(\psi)-\mu_0)
 &=\int_{C^{\An}}t\,\ddc\psi\\
 &=\int_{C^{\An}}\psi\,\ddc t
 =\int_\Gamma\psi|_\Gamma\,\rd\Delta_\Gamma T.
 \end{aligned}
\end{equation}
Only ordinary common model domination is used for the intersection
pairing.  Neither \(\psi\) nor its determining model is required to
factor through a higher reduced semistable skeleton.

For completeness, these rational-breakpoint tests determine a
continuous potential up to a constant.  If a continuous \(w\) satisfies
\(\int w\,\rd\Delta_\Gamma T=0\) for all such \(T\), choose a finite
rational subdivision and let \(T\) interpolate the values of \(w\) at
its vertices.  The measure \(\Delta_\Gamma T\) is supported at those
vertices, so
\[
 0=\int_\Gamma w\,\rd\Delta_\Gamma T
   =\int_\Gamma T\,\rd\Delta_\Gamma T
   =-\sum_{e=(v,v')}\frac{c_e}{\ell_e}
                         \bigl(w(v)-w(v')\bigr)^2.
\]
Connectedness makes all these vertex values equal.  Finer rational
subdivisions have dense union, so continuity makes \(w\) constant
everywhere.  The assertion is immediate for a point skeleton.
Apply this observation to the difference between \(\psi|_\Gamma\)
and the continuous Green potential of
\(p_*\MA_\theta(\psi)-\mu_0\), whose mass is zero.  The kernel of
\cref{lem:ch31-green-kernel} satisfies the same tests, including all
piecewise affine tests.  Thus
\(\Delta_\Gamma(\psi|_\Gamma)=p_*\MA_\theta(\psi)-\mu_0\) in the
stated distributional sense.

\emph{Step 2: the bounded limit.}
Take a decreasing net \(\psi_\alpha\searrow u\) of real
\(\theta\)-psh model functions by
\cref{thm:bfj-decreasing-regularization}.  After restricting to a
tail, this net is uniformly bounded: it lies above the bounded
function \(u\) and below one fixed continuous model function.
The restrictions decrease to \(f=u|_\Gamma\), which is finite and
continuous by clause~(2) of \cref{def:ch21-bfj-psh}.
Dini's theorem for nets, \cref{prop:dini-nets}, gives uniform
convergence on this fixed compact graph.  Meanwhile
\cref{thm:bounded-ma-extension} gives
\(\MA_\theta(\psi_\alpha)\to\MA_\theta(u)\) weakly.
Pushforward by the continuous \(p\) preserves weak convergence,
since a continuous test \(T\) pulls back to \(T\circ p\).

For each fixed rational-breakpoint test \(T\), pass to the limit in
\eqref{eq:ch25-weighted-skeletal-test}.  On the right this is uniform
convergence against the fixed finite signed measure
\(\Delta_\Gamma T\); on the left it is weak convergence.
The result is
\[
 \int_\Gamma T\,\rd(p_*\MA_\theta(u)-\mu_0)
       =\int_\Gamma f\,\rd\Delta_\Gamma T.
\]
The preceding Green-potential and interpolation argument now proves
\(\Delta_\Gamma f=p_*\MA_\theta(u)-\mu_0\).
The measure on the right after adding \(\mu_0\) is positive of mass
\(V\).  The sufficiency direction of
\cref{lem:ch26-weighted-graph-real-bfj} therefore proves that
\(f\circ p=u\circ p\) is continuous BFJ-psh and has measure
\(\iota_*(\mu_0+\Delta_\Gamma f)\).  This is
\eqref{eq:ch25-curve-pushforward}.  The equality of measures, not a
claim about weak convergence on arbitrary Borel sets, gives the
final Borel-set identity.
\end{proof}

\begin{corollary}
\label{cor:ch25-curve-capacity-formula}
Put
\[
 \begin{aligned}
 \mathcal P_\Gamma
  &:=\{u|_\Gamma:u\in\mathcal U_\theta\},\\
 \mathcal Q_\Gamma
  &:=\{f\in C(\Gamma): -1\leq f\leq0,\
                \mu_0+\Delta_\Gamma f\text{ is a positive measure}\}.
 \end{aligned}
\]
Then \(\mathcal P_\Gamma=\mathcal Q_\Gamma\), and for every Borel
\(A\subseteq\Gamma\),
\begin{equation}
\label{eq:ch25-curve-capacity}
 \Capa_\theta(p^{-1}(A))
   =\sup_{f\in\mathcal P_\Gamma}(\mu_0+\Delta_\Gamma f)(A)
   =\sup_{f\in\mathcal Q_\Gamma}(\mu_0+\Delta_\Gamma f)(A).
\end{equation}
Every member of \(\mathcal Q_\Gamma\) pulls back to an actual
continuous BFJ \(\theta\)-psh competitor.  Real values and irrational
breakpoints are allowed.  A real piecewise affine member with
rational breakpoints pulls back to a real model function; if its
vertex values are also rational, it is a rational model function.
\end{corollary}

\begin{proof}
For \(u\in\mathcal U_\theta\),
\cref{prop:ch25-curve-pushforward} gives
\[
 \MA_\theta(u)(p^{-1}(A))
   =(\mu_0+\Delta_\Gamma(u|_\Gamma))(A),
\]
and proves \(\mathcal P_\Gamma\subseteq\mathcal Q_\Gamma\).
Conversely, if \(f\in\mathcal Q_\Gamma\), the graph Laplacian has
total mass zero (test against the constant one), so
\(\mu=\mu_0+\Delta_\Gamma f\) has mass \(V\).
By \cref{lem:ch26-weighted-graph-real-bfj},
\(f\circ p\in\PSH(C,\theta)\cap C^0(C^{\An})\).
It takes values in \([-1,0]\), hence belongs to
\(\mathcal U_\theta\), and restricts to \(f\).
This proves the reverse inclusion and the equality of the suprema.
The final model assertions follow from the SNC construction preceding
\cref{prop:ch25-curve-pushforward}.  No rational approximation of the
real reference class is involved.
\end{proof}

\subsection{The two-vertex graph}
\label{sec:ch25-two-vertex}

Suppose the underlying graph is a segment and \(\mu_0\) is supported
at its two endpoints.  Internal valence-two vertices may be retained
to record changing edge weights, even when their reference masses
are zero.  The coordinate relevant to capacity is the
\emph{resistance coordinate}
\begin{equation}
\label{eq:ch25-resistance-coordinate}
 x(z)=\int_{[x_{E_0},z]}\frac{\rd s}{c_e},\qquad
 \ell=\sum_e\frac{\ell_e}{c_e}.
\end{equation}
Thus \(x\in[0,\ell]\).  Here \(\ell_e\) is an analytic modulus
length, whereas \(\ell\) is a total resistance length; they agree
with ordinary total length only when every \(c_e=1\).
On each edge, if \(f\) is expressed as a function of \(x\), then
\(c_e\,\rd f/\rd s=\rd f/\rd x\).
Accordingly both fluxes at a junction become ordinary one-sided
\(x\)-derivatives, even when its two edge weights differ.

Let \(m_i=\mu_0(\{x_{E_i}\})\).  On the original reduced model these
are the degrees of \(\theta|_{E_i}\) over \(k\); the measure-based
definition also records multiplicities if a different SNC
presentation is used.  Then
\(\mu_0=m_0\delta_0+m_1\delta_\ell\), with \(m_i\geq0\) and
\(V=m_0+m_1>0\).
In the resistance coordinate,
\begin{equation}
\label{eq:ch25-segment-class}
 \mathcal Q_\Gamma
 =\{f\in C([0,\ell]): -1\leq f\leq0,\ f\text{ convex},\
       f'(0^+)\geq-m_0,\ f'(\ell^-)\leq m_1\},
\end{equation}
and
\begin{equation}
\label{eq:ch25-segment-measure}
 \mu_0+\Delta_\Gamma f
  =(m_0+f'(0^+))\delta_0
       +f''|_{(0,\ell)}+(m_1-f'(\ell^-))\delta_\ell .
\end{equation}
Indeed positivity in the interior is precisely convexity in \(x\),
and the remaining two inequalities are positivity at the endpoints.
The endpoint bounds make the one-sided derivatives finite; the
total interior mass is \(f'(\ell^-)-f'(0^+)\), so the displayed
measure has total mass \(V\).

\begin{proposition}
\label{prop:ch25-segment-capacities}
With \(x_c\) denoting the point at resistance distance \(c\) from
\(x_{E_0}\), one has
\begin{align}
\label{eq:ch25-cap-endpoint}
 \Capa_\theta(p^{-1}(\{x_0\}))
   &=m_0+\min\{m_1,1/\ell\},\\
\label{eq:ch25-cap-interval}
 \Capa_\theta(p^{-1}([a,b]))
   &=\min\{m_0,1/a\}+\min\{m_1,1/(\ell-b)\},
       \qquad 0<a\leq b<\ell.
\end{align}
In particular, for \(0<c<\ell\), the capacity of the tube over
\(x_c\) is \(\min\{m_0,1/c\}+\min\{m_1,1/(\ell-c)\}\).
All these suprema are attained by continuous piecewise affine graph
potentials whose pullbacks are BFJ \(\theta\)-psh.
No rationality hypothesis is needed for attainment.
\end{proposition}

\begin{proof}
By \cref{cor:ch25-curve-capacity-formula} it suffices to maximize
the indicated graph mass over \(\mathcal Q_\Gamma\).

For the endpoint, \eqref{eq:ch25-segment-measure} gives the mass
\(m_0+f'(0^+)\).  Convexity implies
\(f'(0^+)\leq f'(\ell^-)\leq m_1\), and
\(\ell f'(0^+)\leq f(\ell)-f(0)\leq1\).
The upper bound in \eqref{eq:ch25-cap-endpoint} follows.
It is attained by \(f(x)=-1+sx\), where
\(s=\min\{m_1,1/\ell\}\): its values lie in \([-1,0]\), and both
endpoint inequalities hold.

For \(0<a\leq b<\ell\), the mass is
\(f''([a,b])=f'(b^+)-f'(a^-)\), including the jumps at \(a,b\).
Convexity and the range bound give
\[
 f'(a^-)\geq\frac{f(a)-f(0)}a\geq-1/a,\qquad
 f'(b^+)\leq\frac{f(\ell)-f(b)}{\ell-b}\leq1/(\ell-b).
\]
The endpoint constraints also give
\(f'(a^-)\geq-m_0\) and \(f'(b^+)\leq m_1\).
Together these prove the desired upper bound.
Set \(\alpha=\min\{m_0,1/a\}\),
\(\beta=\min\{m_1,1/(\ell-b)\}\), and define
\[
 f(x)=
 \begin{cases}
  -1+\alpha(a-x),&0\leq x\leq a,\\
  -1,&a\leq x\leq b,\\
  -1+\beta(x-b),&b\leq x\leq\ell.
 \end{cases}
\]
Its slopes are \(-\alpha,0,\beta\), in increasing order, and its
maximum is attained at an endpoint.  The inequalities
\(\alpha a\leq1\), \(\beta(\ell-b)\leq1\) make that maximum at most
zero; the minimum is minus one.  The endpoint inequalities follow
from \(\alpha\leq m_0,\beta\leq m_1\).
Thus \(f\in\mathcal Q_\Gamma\), and its interior measure is
\(\alpha\delta_a+\beta\delta_b\).
When \(a=b\), these atoms coincide and their masses add.
The corollary makes \(f\circ p\) an admissible competitor even if
\(a,b\) are irrational.  Its measure attains the upper bound
exactly; weak convergence on a fixed Borel set is not needed.
\end{proof}

\begin{remark}[Reading the formulas]
\label{rem:ch25-segment-reading}
The geometry enters through the resistance length, which combines
analytic modulus and residue degree.  For one edge of modulus
length \(L\) and weight \(c\), the term \(1/\ell\) is \(c/L\).
The two terms in \eqref{eq:ch25-cap-interval} are the weighted
fluxes entering from the two ends.  Each minimum records whether
the range constraint or the available reference mass is active.
The formulas are unchanged by subdivisions: lengths of subdivided
pieces add and their inherited weight stays the same.
\end{remark}

\begin{example}[A nonsplit node changes capacity]
\label{ex:ch25-nonsplit-capacity}
Let \(K=\mathbf R((\pi))\), with \(-\log|\pi|=1\), and blow up
\(\mathbf P^1_{\mathbf R[[\pi]]}\) at the closed special-fiber point
\(z\) given by \(x^2+1=0\).
The strict transform \(D\) and the exceptional
\(E\simeq\mathbf P^1_{\mathbf C}\) both have multiplicity one.
They meet at a node with residue field \(\mathbf C\).
The edge therefore has modulus length one and weight two, hence
resistance length \(\ell=1/2\).

Take the pullback of the standard \(\mathcal O(2)\) model as
reference.  Its degrees over \(\mathbf R\) are \((2,0)\), so
\(\mu_0=2\delta_{v_D}\).
The vertical divisor \(-E\) has degree minus two on \(D\) and
plus two on \(E\): the restriction of \(\mathcal O(E)\) to \(D\)
is \(\mathcal O_D(z)\), while its restriction to
\(\mathbf P^1_{\mathbf C}\) is \(\mathcal O(-1)\), of degree
minus two over \(\mathbf R\).
Its model potential has values \((0,-1)\), hence is
\(f(s)=-s=-2x\).  It lies in \([-1,0]\) and has
\[
 \Delta_\Gamma f=-2\delta_{v_D}+2\delta_{v_E},
 \qquad \MA_\theta(f\circ p)=2\delta_{v_E}.
\]
The modified model has degrees \((0,2)\) and is nef.  Therefore
\[
 \Capa_\theta(\{v_E\})
 =\Capa_\theta(p^{-1}(\{v_E\}))=2.
\]
The upper bound is the total mass; the displayed model realizes
it at the point itself.  Equivalently, the endpoint formula from
the \(E\)-end reads \(0+\min\{2,1/(1/2)\}=2\).
Using the modulus length one but omitting its weight would instead
give one.  Thus residue weights affect capacity, not just notation.
\end{example}

